% !TeX root = ../../tesis.tex

% =============================================================================
%  4.3 — VISIÓN GENERAL
% =============================================================================
\section{Transport-gis: El motor de Visualización}
\label{sec:transport-gis-intro}

\textbf{Transport-gis-zmvm-mjg} (
\termino{Transport GIS — Zona Metropolitana del Valle de México, M.J. Garibelt})
es el tercer y último componente de la arquitectura de código de esta
investigación. Mientras \textbf{Apimetro} opera como el motor de datos unificado
—responsable de la ingesta, estandarización y exposición de la red de transporte
como servicio— y \textbf{VFTModel} constituye el núcleo analítico que calcula
los indicadores de desempeño de la red, \textbf{Transport-gis} representa la
capa de visualización: el sistema que traduce los resultados del modelo en
representaciones cartográficas interpretables para el análisis urbanístico.

Al igual que los componentes precedentes, la denominación del sistema proviene
de la obra literaria ficcional \textit{Saga Diarios, tomo~2}, de autoría del
investigador. En dicha novela, la urbanista ficcional
\textit{Mildred Julietta Cordero Garibelt} —referida por su acrónimo
\textit{M.J.G}— es la responsable de la cartografía y la planeación espacial del
sistema de movilidad de \textit{Ciudad Capital}, ucronía de la Ciudad de México
en la década de 1960. La trama sitúa a este personaje en el centro de eventos
políticos relevantes de la historia mexicana del siglo~\textsc{xx}
\citep{lopezmateos2023comoelcantodelcisne}, cuya intervención en las decisiones
sobre infraestructura y movilidad urbana articula el argumento central de la
novela. El paralelismo con la presente investigación es deliberado: así como
\textit{M.J.G} produce los mapas que hacen legible el sistema de movilidad de
\textit{Ciudad Capital}, \textbf{Transport-gis} genera las representaciones
cartográficas que permiten interpretar los resultados del \textbf{VFTModel}
sobre la red metropolitana real de la \zmvm.
